% Page budget: approximately 4 pages.
\section{Dynamic library model and sufficient state}\label{sec:model}

The model is specified first on raw verified libraries.  Research then
generates and verifies a catalog candidate, admission updates the raw library,
and the compressed process is obtained by projecting that raw evolution.

\subsection{Hidden market state and filtered belief}

Let \(X=\Fin(m)\), \(2\le m\), be the hidden market-state type and let \(B\) be
a nonempty finite grid of filtered beliefs.  The controller observes \(b\in B\)
rather than \(x\in X\).  Each grid point has an injective rational
interpretation \(b\mapsto\mu_b\), at least one \(\mu_b\) is non-Dirac, and the
already-filtered belief evolves under an exact rational Markov kernel \(P_B\).
Economically, \(x\) is the unobserved market regime and \(b\) is the
controller's decision-relevant assessment of that regime.  All carriers below
are finite and have decidable equality; the probability-table and
normalization conventions are in \cref{app:ratprob}.

\subsection{Strategy operating profiles}

Each strategy identifier \(s\in S\) has an immutable state-contingent operating
payoff \(u_s:X\to\Q\).  Its belief-indexed \emph{operational profile} is
\begin{equation}\label{eq:operational-profile}
  j_s(b):=\sum_{x\in X}\mu_b(x)u_s(x).
\end{equation}
Thus \(j_s(b)\) is the expected payoff from operating \(s\) now; project and
continuation values are introduced separately below.  A distinguished
inactive strategy \(s_0\) has zero operating payoff.

\subsection{Modules and closure}

Each strategy also carries an immutable finite set of modules, so its complete
catalog row is
\[
  \sigma_s=(s,u_s,\mods(s)),
  \qquad \mods(s)\in\powersetfin(M),
  \qquad \mods(s_0)=\varnothing.
\]
Economically, modules are reusable capabilities: they do not enter operating
payoff directly, but they can be combined to support later research.  The set
of capabilities available after all such combinations is represented by an
extensive, monotone, and idempotent closure operator
\[
  \cl:\powersetfin(M)\to\powersetfin(M).
\]

\subsection{Raw library}

A \emph{raw verified library} is a set \(L\subseteq S\) containing \(s_0\).
Economically, \(L\) is the catalog of strategy identities currently retained
and verified.  Define
\begin{align}
  F_L(b)&:=\max_{s\in L}j_s(b),
  &U_L&:=\bigcup_{s\in L}\mods(s),\label{eq:frontier-union}\\
  C_L&:=\cl(U_L).\label{eq:raw-library-closure}
\end{align}
Here \(F_L\) is the best currently attainable operating profile, \(U_L\) is
the raw module inventory, and \(C_L\) is the resulting set of generative
capabilities.  The raw decision state is \((b,L)\).  Continue earns \(F_L(b)\),
advances the belief once under \(P_B\), and leaves \(L\) unchanged.  The
elementary frontier and closure calculus is recorded in
\cref{lem:frontier-closure-calculus}.

\subsection{Positive-duration generation, verification, and admission}

Economically, a research project uses the incumbent library's capabilities
over a positive number of periods to seek a new verified strategy.  Each
project \(q\in Q\) has prerequisites
\(\operatorname{req}(q)\in\powersetfin(M)\).  At initiation belief \(b\) and
closure \(C\), candidate generation follows
\[
  G(q,b,C)\in\Delta_{\Q}(\operatorname{Option}(S)),
\]
with the point mass on \(\operatorname{none}\) when the prerequisites are not
contained in \(C\).  A generated \(s\) passes verification with probability
\(\nu(q,b,C,s)\).  The admitted law is derived from these two stages:
\[
  \Gamma(q,b,C)(\operatorname{some}(s))
  =G(q,b,C)(\operatorname{some}(s))\nu(q,b,C,s),
\]
with all failed-generation and failed-verification mass assigned to
\(\operatorname{none}\).  Because \(G\) sees only \((q,b,C)\), generation cannot
inspect raw identifiers, multiplicity, provenance, age, or admission order.

For a frontier--closure pair \(K=(F,C)\), let \(A(K)\) be the available finite
project menu.  A project \(q\in A(K)\) has initiation cost
\(\kappa_q(b,K)\ge0\), strictly positive duration \(d_q\), and an operation
flag \(o_q\in\{0,1\}\).  Its exact joint completion law
\(\Lambda_q(\mathbf b,o\mid b,K)\) couples the \(P_B\)-belief path with the
admitted outcome \(o\sim\Gamma(q,b,C)\); the components may be correlated.
During research the incumbent earns
\[
  -\kappa_q(b,K)+\sum_{t=0}^{d_q-1}\beta^t o_qF(b_t),
  \qquad 0\le\beta<1.
\]
The old library remains in force through date \(d_q-1\), the outcome is
applied at date \(d_q\), and continuation starts from the terminal belief.
Appendix~A gives the full path marginals and normalization identities.

\subsection{Raw update}

Economically, successful admission adds one verified candidate, whereas every
other outcome leaves the incumbent catalog unchanged.  Formally,
\begin{equation}\label{eq:raw-library-update}
  L\oplus\operatorname{none}=L,
  \qquad
  L\oplus\operatorname{some}(s)=L\cup\{s\}.
\end{equation}
The corresponding raw terminal law is the pushforward
\begin{equation}\label{eq:raw-terminal-transition}
  \mathcal Q_q^{\rm raw}(\,\cdot\mid b,L)
  :=(\mathbf b,o\mapsto(b_{d_q},L\oplus o))_{\#}
    \Lambda_q(\,\cdot\mid b,(F_L,C_L)).
\end{equation}

\subsection{Compressed frontier--closure state}

\begin{definition}[Compressed frontier--closure state]
\label{def:finite-library-state}
The compressed innovation state is
\begin{equation}\label{eq:closure-compression}
  K_L:=(F_L,C_L).
\end{equation}
Economically, the frontier records what the library can operate now and the
closure records what it can help generate later.  The realizable set
\(\mathcal K:=\{K_L:L\text{ is a raw verified library}\}\) is finite because
\(S\) is finite.
\end{definition}

For an admitted outcome define
\[
  \operatorname{addK}(K,\operatorname{none})=K,
  \qquad
  \operatorname{addK}((F,C),\operatorname{some}(s))
  =\left(b\mapsto\max\{F(b),j_s(b)\},
    \cl(C\cup\mods(s))\right).
\]
This map is the compressed effect of adding an admitted catalog row.  Closure
idempotence gives the local identity
\begin{equation}\label{eq:raw-compressed-local-update}
  K_{L\oplus o}=\operatorname{addK}(K_L,o).
\end{equation}
Applying \(L\mapsto K_L\) to the raw law in
\cref{eq:raw-terminal-transition} yields the next-state laws
\(\overline T_q\) and \(\overline{\mathcal Q}_q\): they are the pushforwards of
\(\Gamma\) and \(\Lambda_q\) through \(\operatorname{addK}\) and
\((\mathbf b,o)\mapsto(b_{d_q},\operatorname{addK}(K,o))\), respectively.
Thus the compressed transition is derived from the raw law, not separately
postulated.

\subsection{Projection theorem}

Economically, two raw libraries with the same frontier--closure state support
the same operating and research opportunities even when their retained
strategy identities differ.

\begin{theorem}[Raw-to-compressed controlled Markov projection]
\label{thm:raw-to-compressed-projection}
In the finite exact model above:
\begin{enumerate}
  \item \(\overline T_q\) and \(\overline{\mathcal Q}_q\) are normalized and
    supported on realizable compressed states;
  \item if \(K_L=K_{\widetilde L}\), every project has the same availability,
    cost, duration, operating-reward block, and joint terminal
    belief--next-compressed-state law from \(L\) and \(\widetilde L\);
  \item at embedded decision epochs,
    \((B_{\tau_n},K_{L_n})\) is the controlled semi-Markov projection of
    \((B_{\tau_n},L_n)\); and
  \item for every finite calendar horizon \(h\), belief \(b\), and raw library
    \(L\),
    \[
      V_h^{\rm raw}(b,L)=\bar V_h(b,K_L).
    \]
    The stationary raw and compressed Bellman fixed points also agree under
    \(L\mapsto K_L\), and a maximizing stationary compressed selector lifts to
    a Bellman-optimal raw selector.
\end{enumerate}
\end{theorem}

The projection, quotient, and restricted representation-refinement arguments
are in \cref{app:di-quotient}.  Finiteness supplies normalization
and attained action maxima, positive duration makes the calendar-horizon
induction well founded, and contraction is used only for the stationary
clauses.  The proof leaves the joint completion law unfactorized, so it permits
correlation between the belief path and admission outcome.  Detectability does
not enter this result.  The Markov state is stated at decision epochs; during
an unfinished project it must also record the project and remaining duration.

\subsection{Resource burden and exact compression}
\label{subsec:retention-resources}

Resource burden is an outer retention decision, not operating payoff or
project cost.  Each active strategy has a fixed rational weight \(w_s>0\),
while \(w_{s_0}=0\), and the library's one-time retention burden is
\begin{equation}\label{eq:library-burden}
  W(L):=\sum_{s\in L}w_s.
\end{equation}
Economically, \(W(L)\) can commensurate storage, validation, maintenance,
monitoring, retrieval, and governance requirements.  If
\(S_\theta^{\mathrm{elig}}\) is the finite outer-certified catalog, write
\(\mathfrak L_\theta=\{L'\subseteq S_\theta^{\mathrm{elig}}:s_0\in L'\}\).
Later capacity and price problems optimize over this family.  Equal compressed
states can have unequal burdens, so a resource optimizer must retain the raw
library \(L\) or augment \(K_L\) by \(W(L)\).

For compression of a specified source \(L\), define
\begin{equation}\label{eq:safe-compression-family}
  \mathcal P_{\mathrm{safe}}(L)
  :=\{L'\subseteq L:s_0\in L',\ K_{L'}=K_L\}.
\end{equation}
This family contains exactly the source sublibraries that preserve both the
current operating frontier and the generative closure.  By the projection
theorem, every member therefore has the same project menus,
positive-duration research laws, future compressed updates, and productive
values as the source.  Resource minimization changes only the retained raw
representation of that control process.  The detectability result in
\cref{sec:compression} gives the corresponding observable converse.  Exact
safe compression is
\begin{equation}\label{eq:minimum-safe-compression}
  \boxed{
  \begin{aligned}
    \underset{L'\subseteq L}{\operatorname{minimize}}\quad &W(L')\\
    \text{subject to}\quad &s_0\in L',\\
                            &K_{L'}=K_L.
  \end{aligned}}
\end{equation}
